\documentclass[11pt,a4paper]{article}

\usepackage[T1]{fontenc}
\usepackage[utf8]{inputenc}
\usepackage{lmodern}
\usepackage[french]{babel}
\usepackage[a4paper,margin=2.3cm]{geometry}
\usepackage{microtype}
\usepackage{xcolor}
\usepackage{enumitem}
\usepackage{titlesec}
\usepackage{hyperref}
\usepackage{tabularx}
\usepackage{array}
\usepackage{booktabs}
\usepackage{fancyvrb}
\usepackage[most]{tcolorbox}

\definecolor{ink}{HTML}{1c2330}
\definecolor{accent}{HTML}{2f6e4c}
\definecolor{accent2}{HTML}{8a5a1c}
\definecolor{soft}{HTML}{eef3ef}
\definecolor{softline}{HTML}{cfe0d4}
\definecolor{open}{HTML}{1c4e7a}
\definecolor{opensoft}{HTML}{eaf1f8}
\definecolor{theme}{HTML}{5a2f6e}
\definecolor{themesoft}{HTML}{f2ecf6}
\definecolor{codebg}{HTML}{f4f6f4}
\definecolor{codeline}{HTML}{d6e2d8}

\hypersetup{colorlinks=true, linkcolor=accent, urlcolor=accent2,
  pdftitle={Concept --- Acte 1 / Prologue}}

\titleformat{\section}{\Large\bfseries\color{ink}}{\thesection.}{0.5em}{}
\titleformat{\subsection}{\large\bfseries\color{accent}}{\thesubsection}{0.6em}{}
\titlespacing{\section}{0pt}{1.3em}{0.45em}
\setlist[itemize]{leftmargin=1.2em,itemsep=2pt,topsep=2pt}

\newtcolorbox{pitch}{enhanced,breakable,colback=soft,colframe=accent,
  boxrule=0.8pt,left=10pt,right=10pt,top=7pt,bottom=7pt,arc=3pt}
\newtcolorbox{defer}{enhanced,breakable,colback=opensoft,colframe=open,
  boxrule=0.5pt,left=9pt,right=9pt,top=6pt,bottom=6pt,arc=2pt}

% bloc de code / donnée
\DefineVerbatimEnvironment{code}{Verbatim}{fontsize=\small,frame=single,
  framerule=0.5pt,rulecolor=\color{codeline},formatcom=\color{ink}}

\title{\vspace{-1.5em}\textbf{\color{ink}\Huge Concept --- Acte 1 / Prologue}\\[0.3em]
\large\color{accent}Première brique jouable de la campagne\\[0.2em]
\normalsize\color{ink}\itshape Plan de construction --- traduit la direction et le thème en structure concrète}
\author{}
\date{}

\begin{document}
\maketitle
\vspace{-3em}

\begin{pitch}
\textbf{\color{ink}Pitch.} On appartient à un \textbf{camp} dont dépend un petit \textbf{territoire}, découpé en \textbf{${\sim}$12 régions} (cellules de maillage) : c'est le \textbf{geoscape de départ}. L'acte ouvre sur une \textbf{mission d'intro} (elle pose les personnages et le ton --- la réalité se déforme, cf. \texttt{Opening.txt}), puis le joueur arrive sur le \textbf{geoscape}, voit sa carte, et \textbf{choisit} par où continuer.
\end{pitch}

\noindent\textbf{La boucle de l'acte :}
\begin{code}
Intro (mission) -> GEOSCAPE (voir la carte, choisir) -> Mission de region
                       ^                                      |
                       +---------- retour (win ET lose) <-----+
   ... jusqu'a la condition de fin d'acte  ->  Acte 2 / interlude
\end{code}

% ======================================================================
\section{Le territoire (la carte geoscape)}
\begin{itemize}
\item \textbf{${\sim}$12 cellules.} Une = \textbf{le camp} (le hub dont dépend le territoire, point de retour). Les autres = \textbf{régions}.
\item Chaque cellule porte : \texttt{nom}, \texttt{type} (camp | région), \texttt{contenu} (mission liée | événement | vide), \texttt{état} (verrouillé | disponible | nettoyé | perdu/menacé), et ses \texttt{voisins} (= l'adjacence naturelle du maillage).
\item \textbf{Géométrie :} maillage \textbf{grossier généré} (seed) via \texttt{genMesh}, puis cellules \textbf{nommées/annotées à la main}.
\item \textbf{Choix réel dès l'intro} (pilier A3) : à la sortie de l'intro, \textbf{$\geq$ 2-3 régions disponibles} simultanément --- sinon «~choisir~» est factice.
\end{itemize}

% ======================================================================
\section{Les pièces à construire (architecture décidée)}

\subsection{A --- Format de donnée geoscape \textnormal{\small(nouveau --- dossier \texttt{geoscapes-mesh/})}}
Le maillage réutilise le \texttt{form} des missions (\texttt{infl}/\texttt{dens}/\texttt{seed}) ; on annote les cellules :
\begin{code}
{
  "name": "Acte 1 - Territoire du camp",
  "mesh": { "form": { "infl": [...], "dens": 90, "seed": 1234 } },
  "camp": "c0",
  "cells": [
    { "id": "c0", "name": "Le Camp", "type": "camp",
      "content": { "kind": "empty" }, "state": "available" },
    { "id": "c3", "name": "Bois brule", "type": "region",
      "content": { "kind": "mission", "ref": "mission-1.json" },
      "state": "available", "unlockedBy": [] },
    { "id": "c7", "name": "Gue de l'aval", "type": "region",
      "content": { "kind": "mission", "ref": "mission-2.json" },
      "state": "locked", "unlockedBy": ["c3"] }
  ]
}
\end{code}
\emph{(les voisins sont dérivés de la géométrie, pas stockés à la main)}

\subsection{B --- Nœud \texttt{geoscape} dans le graphe de campagne \textnormal{\small(extension de l'éditeur)}}
Un nouveau \texttt{type} de nœud, à côté de \texttt{text} et \texttt{mission} :
\begin{code}
{ "id": "g1", "type": "geoscape", "map": "acte1.json",
  "endWhen": { "cleared": ["c11"] },   // condition de fin d'acte
  "next": "n_acte2",                   // ou aller quand endWhen est vrai
  "x": ..., "y": ... }
\end{code}
\begin{itemize}
\item L'intro \texttt{win $\to$ g1}.
\item À l'entrée de \texttt{g1} : on \textbf{affiche la carte}. Le joueur clique une région \texttt{disponible} $\to$ lance la mission liée (\texttt{content.ref}).
\item La mission de région \texttt{win/lose $\to$} \textbf{revient à \texttt{g1}} (et met à jour les états : \texttt{cleared}, déverrouillage des voisins\dots).
\item \textbf{Échec qui fait avancer (C4) :} \texttt{lose} revient aussi à \texttt{g1}, mais marque la région \texttt{perdu}/monte une menace --- \textbf{jamais} un reload.
\item Quand \texttt{endWhen} est satisfait $\to$ \texttt{next} (Acte 2 / interlude).
\end{itemize}

\subsection{C --- Runtime geoscape jouable \textnormal{\small(le livrable «~on voit la map~»)}}
L'écran qui : charge la map, \textbf{rend le maillage}, colorie les états (dispo / verrouillé / nettoyé / perdu), montre les régions cliquables + leur info (nom, type de mission, menace), lance la mission, et y revient.

\subsection{D --- Persistance du roster}
Les \textbf{mêmes personnages nommés} d'une mission à l'autre. La \emph{progression} (perks) est différée, mais \textbf{l'identité et le roster persistent} dans l'état de campagne (indispensable pour C1/B9 : on s'attache, on perd).

% ======================================================================
\section{Ancrage direction --- les piliers actifs dès l'Acte 1}
\renewcommand{\arraystretch}{1.3}
\begin{tabularx}{\linewidth}{@{}>{\bfseries}p{4.2cm} X@{}}
\toprule
\textmd{\textbf{Pilier}} & \textbf{Présence en Acte 1} \\
\midrule
A3 Choix concurrent & $\geq$ 3 régions disponibles à la sortie de l'intro \\
C4 Échec fait avancer & \texttt{lose} $\to$ état changé (région perdue / menace), pas de reload \\
C1/B9 Identité \& mort & Roster nommé persistant ; statut spécial des persos clés de l'intro \\
Thème (déformation) & La carte du monde \textbf{est} un maillage déformé ; (option) elle se déforme au fil de l'acte \\
A1 Horloge locale & \emph{(option, différable)} une menace par région qui monte si on tarde \\
\bottomrule
\end{tabularx}
\renewcommand{\arraystretch}{1}

% ======================================================================
\section{Contenu minimal pour la tranche Acte 1}
\begin{itemize}
\item \textbf{1 camp} + \textbf{intro} --- déjà en place : \texttt{Opening.txt}, \texttt{mission-1.json}, \texttt{Mission1-beginning/end.txt}.
\item \textbf{3 régions} avec mission (dont au moins une réutilise \texttt{mission-1} au départ).
\item \textbf{1 région-clé} dont le nettoyage = \textbf{fin d'acte}.
\item \textbf{Interludes} déjà écrits (\texttt{Interlude1\dots4}) branchés aux transitions.
\end{itemize}

% ======================================================================
\section{Ordre de construction (après ce doc)}
\begin{enumerate}
\item \textbf{Format geoscape + éditeur dédié} : génération grossière (\texttt{genMesh}), nommage des cellules, marquage du camp, liaison cellule$\to$mission, états/déverrouillages. \emph{(la fondation)}
\item \textbf{Nœud \texttt{geoscape}} dans l'éditeur de campagne + câblage du retour (win/lose $\to$ geoscape) et de \texttt{endWhen}.
\item \textbf{Runtime geoscape jouable} : rendu + sélection + lancement + retour.
\item \textbf{Câbler le contenu minimal} de l'Acte 1 et \textbf{tester la boucle complète} (intro $\to$ carte $\to$ choix $\to$ mission $\to$ retour $\to$ fin d'acte).
\end{enumerate}

% ======================================================================
\section{Différé --- à décider plus tard}
\begin{defer}
\begin{itemize}
\item \textbf{Progression des personnages} (perks A/B par grade) --- explicitement remise à après cette boucle.
\item \textbf{Détail de l'horloge locale} (A1) --- peut rester minimal en Acte 1.
\item \textbf{Déformation visuelle progressive} de la carte (levier thème) --- à prototyper une fois la boucle debout.
\item \textbf{Dosage du hasard en combat} (friction 1, ouverte) --- tranchée au prototype combat, indépendante de cet acte.
\end{itemize}
\end{defer}

\end{document}
